\section*{Bab \ref{ch:g11:subduction-zones} --- Zona Subduksi}

\begin{solution}{exo:g11:subduction-zones:1}
\omterm{def:g11:subduction-zones:subduction}{Subduksi} adalah tenggelamnya lempeng \omterm{def:g11:ocean-ridges:oceanic-lithosphere}{litosfer samudra} di bawah lempeng lain,
kembali ke dalam mantel; \omterm{def:g11:subduction-zones:subduction}{zona subduksi} adalah batas lempeng konvergen tempat
\omterm{def:g11:subduction-zones:subduction}{subduksi} terjadi, yang ditandai oleh palung. \omterm{def:g11:ocean-ridges:oceanic-lithosphere}{Litosfer samudra} yang tenggelam:
setelah tua, \omterm{def:g11:lithosphere-asthenosphere:lithosphere}{litosfer} itu dingin, tebal, dan lebih rapat daripada \omterm{def:g11:lithosphere-asthenosphere:lithosphere}{astenosfer} di
bawahnya. Kerak benua terlalu ringan untuk tenggelam.
\end{solution}

\begin{solution}{exo:g11:subduction-zones:2}
(1) Misalnya wilayah Tonga--Fiji, Pasifik barat dari Jepang sampai Kepulauan
Mariana, Indonesia, dan Amerika Selatan. (2) Selalu di sisi lempeng yang berada
di atas, tidak pernah di sisi samudra. (3) Dangkal: semua gempa bumi di Punggung
Tengah Atlantik pada peta itu lebih dangkal daripada \qty{70}{km}.
\end{solution}

\begin{solution}{exo:g11:subduction-zones:3}
$7031 \times 0.21 \approx 1480$ gempa bumi lebih dalam daripada \qty{70}{km};
$7031 \times 0.069 \approx 485$ lebih dalam daripada \qty{300}{km}.
\end{solution}

\begin{solution}{exo:g11:subduction-zones:4}
\omterm{def:g11:subduction-zones:eclogite}{Sekis biru} adalah batuan metamorf yang terbentuk dari basalt \omterm{def:g11:ocean-ridges:oceanic-lithosphere}{kerak samudra} yang
terkubur sampai bertekanan tinggi sementara tetap relatif dingin, yang diwarnai
biru oleh glaukofan; \omterm{def:g11:subduction-zones:eclogite}{eklogit} adalah batuan metamorf yang rapat yang terbentuk
dari basalt dan \omterm{def:g11:ocean-ridges:gabbro}{gabro} pada kedalaman mantel, yang terdiri atas garnet merah dan
piroksen hijau. Keduanya berasal dari \omterm{def:g11:ocean-ridges:oceanic-lithosphere}{kerak samudra} dan mempertahankan komposisi
kimianya; yang berbeda adalah mineral-mineralnya, yang terbentuk dalam keadaan
padat di bawah tekanan dan suhu yang baru, serta massa jenisnya.
\end{solution}

\begin{solution}{exo:g11:subduction-zones:5}
Dasar samudra: \omterm{def:g11:lithosphere-asthenosphere:lithosphere}{litosfer} tua yang bergerak ke arah palung. Palung: lempeng
melengkung turun dan mulai tenggelam. Busur muka: sedimen yang terkikis dari
lempeng yang tenggelam menumpuk; di bawahnya, kedua lempeng saling bergeser pada
batas lempeng, tempat terjadinya gempa-gempa bumi besar. Busur gunung api:
deretan gunung api di atas \omterm{def:g11:subduction-zones:slab}{lempengan}. Busur belakang: di belakang sebagian busur,
lempeng yang berada di atas teregang.
\end{solution}

\begin{solution}{exo:g11:subduction-zones:6}
$\tan\delta = (508 - 135)/(800 - 300) = 373/500 \approx 0.746$, jadi
$\delta \approx \ang{37}$.
\end{solution}

\begin{solution}{exo:g11:subduction-zones:7}
(1) $\tan\ang{45} = 1$: kedalaman \qty{300}{km}. (2) Pada jarak \qty{670}{km}
dari palung. (3) \omterm{def:g11:subduction-zones:slab}{Lempengan} sepanjang $670\sqrt{2} \approx \qty{948}{km}$.
\end{solution}

\begin{solution}{exo:g11:subduction-zones:8}
(1) $(74.6 + 70.0 + 65.5 + 63.1 + 52.6)/5 \approx \qty{65.2}{mm}$ setahun;
$(47.6 + 31.0 + 29.2 + 22.8 + 19.8 + 15.2 + 14.9 + 10.6)/8 \approx
\qty{23.9}{mm}$ setahun; perbandingannya $65.2/23.9 \approx 2.7$.
(2) Lempeng-lempeng bergerak paling cepat bila padanya tergantung \omterm{def:g11:subduction-zones:slab}{lempengan}:
\omterm{def:g11:subduction-zones:slab-pull}{tarikan lempengan} yang tenggelam adalah gaya terkuat yang menggerakkan lempeng,
seperti yang disimpulkan Forsyth dan Uyeda.
\end{solution}

\begin{solution}{exo:g11:subduction-zones:9}
(1) \qty{83}{mm} sama dengan \qty{0.083}{m}: $50/0.083 \approx 600$ tahun dan
$60/0.083 \approx 720$ tahun gerak lempeng. (2) Dari 869 sampai 2011 ada 1142
tahun: ordenya sama. Jika batas lempeng tetap terkunci di antara gempa-gempa bumi
besar, pergeseran tahun 2011 melepaskan gerak yang tersimpan selama beberapa
abad.
\end{solution}

\begin{solution}{exo:g11:subduction-zones:10}
(1) Jepang: $6400 - 3200\,e^{-133/62.8} \approx \qty{6015}{m}$; Tonga:
$6400 - 3200\,e^{-99/62.8} \approx \qty{5740}{m}$; Andes:
$2500 + 350\sqrt{49} = \qty{4950}{m}$. (2) Dasar yang mencapai Jepang: yang
paling tua, jadi \omterm{def:g11:lithosphere-asthenosphere:lithosphere}{litosfer} yang paling dingin dan paling rapat, serta dasar yang
paling dalam.
\end{solution}

\begin{solution}{exo:g11:subduction-zones:11}
(1) $3380 \times 1.06 \approx \qty{3583}{kg/m^3}$. (2)
$(3583 - 2900)/2900 \approx 0.24$: kenaikan sekitar \qty{24}{\percent}.
(3) Sebelumnya, kerak lebih ringan daripada mantel sebesar
$3380 - 2900 = \qty{480}{kg/m^3}$; sesudahnya, kerak lebih rapat sekitar
\qty{203}{kg/m^3}. Kerak itu tidak lagi menahan \omterm{def:g11:subduction-zones:slab}{lempengan}: kerak itu menariknya
ke bawah.
\end{solution}

\begin{solution}{exo:g11:subduction-zones:12}
(1) $t = 6 \times 10^6 \times 3.16 \times 10^7 \approx 1.9 \times 10^{14}$~s;
$\sqrt{10^{-6} \times 1.9 \times 10^{14}} \approx \qty{1.4e4}{m}$, sekitar
\qty{14}{km}. (2) Jauh lebih kecil daripada \qty{100}{km}: dalam 6 juta tahun
hanya kulit luar \omterm{def:g11:subduction-zones:slab}{lempengan} yang menghangat; bagian dalamnya tetap dingin.
(3) Batuan yang dingin itu kaku dan dapat patah, sedangkan mantel yang panas di
sekitarnya mengalir; gempa-gempa bumi yang dalam terjadi di dalam \omterm{def:g11:subduction-zones:slab}{lempengan} yang
dingin itu.
\end{solution}

\begin{solution}{exo:g11:subduction-zones:13}
(1) $\sqrt{1017^2 + 670^2} \approx \qty{1218}{km}$. (2) \qty{64}{mm} setahun
sama dengan \qty{64}{km} per sejuta tahun: $1218/64 \approx 19$ juta tahun.
(3) Sekitar tiga kali lebih lama daripada di bawah Tonga: \omterm{def:g11:subduction-zones:slab}{lempengan} di bawah
Andes tenggelam lebih lambat dan dengan kemiringan yang lebih landai.
\end{solution}

\begin{solution}{exo:g11:subduction-zones:14}
(1) Irisan selebar \qty{1}{m} di sepanjang palung mempunyai volume
$10^5 \times 10^6 \times 1 = \qty{1e11}{m^3}$; kelebihan massanya
$36 \times 10^{11} = \qty{3.6e12}{kg}$. (2) Kelebihan beratnya
$3.6 \times 10^{12} \times 9.8 \approx \qty{3.5e13}{N}$ untuk setiap meter
palung. (3) \omterm{def:g11:subduction-zones:slab}{Lempengan} itu terikat pada bagian lain lempeng, yang kaku: gaya itu
diteruskan ke seluruh lempeng, yang terseret ke arah palung, terlebih lagi karena
palung-palung membentang ribuan kilometer.
\end{solution}

\begin{solution}{exo:g11:subduction-zones:15}
(1) $(23.6 - 16.5) \times 111 \approx \qty{790}{km}$. (2) Laju itu bertambah
$237 - 104 = \qty{133}{mm}$ setahun sepanjang \qty{790}{km}, sekitar
\qty{17}{mm} setahun per seratus kilometer. (3) Laju saling mendekat adalah
kecepatan lempeng Pasifik dikurangi kecepatan busur. Karena kecepatan lempeng
Pasifik kira-kira sama di kedua titik, busur itu sendiri pasti bergerak ke arah
lempeng Pasifik, makin cepat ke utara: busur belakang di belakangnya sedang
teregang, lebih banyak di utara daripada di selatan.
\end{solution}

\begin{solution}{pb:g11:subduction-zones:1}
\textbf{1.} Sekitar \qty{671}{km}.

\textbf{2.} Di sisi samudra lempeng itu belum tenggelam: \omterm{def:g11:lithosphere-asthenosphere:lithosphere}{litosfer} di sana berada
di dekat permukaan. Gempa-gempa bumi yang dalam terjadi di dalam \omterm{def:g11:subduction-zones:slab}{lempengan}, yang
hanya terletak di bawah lempeng yang berada di atas.

\textbf{3.} $z(300) = -24 + 273 = \qty{249}{km}$ dan
$z(700) = -24 + 637 = \qty{613}{km}$; $\tan\delta = 364/400 = 0.91$, jadi
$\delta \approx \ang{42}$.

\textbf{4.} $x = (100 + 24)/0.91 \approx \qty{136}{km}$ dan
$x = (400 + 24)/0.91 \approx \qty{466}{km}$.

\textbf{5.} Garis lurus dengan kemiringan ini yang bermula di palung akan melalui
$z = 0$ pada $x = 0$; garis penampang itu terletak di atasnya. \omterm{def:g11:subduction-zones:slab}{Lempengan} itu
tidak lurus: di dekat palung kemiringannya lebih landai, dan makin curam di
bagian yang lebih dalam.

\textbf{6.} \qty{16.8}{cm} setahun, \qty{168}{km} per sejuta tahun.

\textbf{7.} $168/69 \approx 2.4$ kali lebih cepat.

\textbf{8.} Laju saling mendekat melintasi palung adalah kecepatan lempeng
Pasifik dikurangi kecepatan busur di sisi seberangnya. Jika busur itu bergerak ke
arah lempeng Pasifik, seperti yang terjadi selama busur belakang di belakangnya
teregang, kedua kecepatan itu saling menjumlah, dan laju saling mendekat lebih
besar daripada kecepatan lempeng Pasifik saja.

\textbf{9.} $6400 - 3200\,e^{-99/62.8} \approx \qty{5740}{m}$.

\textbf{10.} Setelah sekitar 60 juta tahun \omterm{def:g11:ocean-ridges:oceanic-lithosphere}{litosfer samudra}, secara keseluruhan,
lebih rapat daripada \omterm{def:g11:lithosphere-asthenosphere:lithosphere}{astenosfer} di bawahnya; pada umur 99 juta tahun \omterm{def:g11:lithosphere-asthenosphere:lithosphere}{litosfer} itu
lebih rapat lagi. Begitu tepinya turun, beratnya menarik bagian selebihnya masuk.

\textbf{11.} $(670 + 24)/0.91 \approx \qty{763}{km}$.

\textbf{12.} $\sqrt{763^2 + 670^2} \approx \qty{1015}{km}$.

\textbf{13.} $1015/168 \approx 6.0$ juta tahun.

\textbf{14.} $1015/69 \approx 14.7$ juta tahun.

\textbf{15.} $3380 \times 1.06 \approx \qty{3583}{kg/m^3}$: kerak, yang semula
lebih ringan daripada mantel, menjadi lebih rapat daripadanya, dan seluruh
\omterm{def:g11:subduction-zones:slab}{lempengan} menarik lebih kuat.

\textbf{16.} $t \approx 6.0 \times 10^6 \times 3.16 \times 10^7 \approx
1.9 \times 10^{14}$~s, dan $\sqrt{\kappa t} \approx \qty{14}{km}$.

\textbf{17.} \qty{14}{km} itu kecil dibandingkan dengan \qty{100}{km}: hanya
bagian luar \omterm{def:g11:subduction-zones:slab}{lempengan} yang telah menghangat; bagian dalamnya masih dingin ketika
mencapai \qty{670}{km}.

\textbf{18.} \omterm{def:g11:subduction-zones:slab}{Lempengan} yang dingin tetap cukup kaku untuk patah, sampai ke
ujungnya, sementara mantel yang panas di sekitarnya mengalir: gempa-gempa bumi di
bawah Tonga terjadi di dalam \omterm{def:g11:subduction-zones:slab}{lempengan}, sampai kedalaman sekitar \qty{671}{km}.

\textbf{19.} Gelombang seismik menunjukkan perubahan yang tajam di mantel pada
sekitar \qty{670}{km}: kecepatannya melonjak di sana.

\textbf{20.} Sekitar \textbf{6 juta tahun} dari palung sampai
\qty{670}{km}.
\end{solution}
